% PAGES: 261-266
% This is the whole of Chapter 8's exercise list (items 1-8, PDF pages
% 261-265 / folios 245-249), kept in a single \exercises environment so the
% shared enumerate counter is never reset mid-list. PDF page 266 (folio 250)
% is genuinely blank: rendered and inspected directly at 300 dpi, no text or
% rule of any kind, only scanner speckle. PDF page 267 (folio 251) was also
% rendered and confirmed to open "Chapter 9" fresh -- this is the true end of
% Chapter 8, not a dropped page. This is the last chunk of the chapter: no
% environment is left open at the end of this file, and nothing further
% follows.

\section{Exercises}

\begin{exercises}

\item Denote by $x^*$ the solution to the ordinary least squares problem $y\approx Ax$,
where $A$ is an $m\times n$ matrix with $m>n$, $y$ is an $m\times1$ column vector, and
$x$ is an $n\times1$ column vector.

Show that $x^*$ is the projection of the vector $y$ onto the space generated by the
columns of $A$, or, equivalently, show that the vector $y-Ax^*$ is orthogonal to the
space generated by the columns of $A$. In other words, show that
\[
(y-Ax^*, Az) \;=\; 0, \quad \forall\, z\in\R^n.
\]

Note: Recall that the space generated by the columns of $A$ is $\{Az \mid z\in\R^n\}$,
and is also called the range of the matrix $A$.

\item Recall from the example from the book that, on March 9, 2012, the mid prices of
the S\&P 500 options maturing on December 22, 2012, were as follows:
\begin{center}
\begin{tabular}{ccc}
Call Price & Strike & Put Price \\[0.4em]
225.40 & 1175 & 46.60 \\
205.55 & 1200 & 51.55 \\
186.20 & 1225 & 57.15 \\
167.50 & 1250 & 63.30 \\
149.15 & 1275 & 70.15 \\
131.70 & 1300 & 77.70 \\
115.25 & 1325 & 86.20 \\
99.55 & 1350 & 95.30 \\
84.90 & 1375 & 105.30 \\
71.10 & 1400 & 116.55 \\
58.70 & 1425 & 129.00 \\
47.25 & 1450 & 143.20 \\
29.25 & 1500 & 173.95 \\
15.80 & 1550 & 210.80 \\
11.10 & 1575 & 230.90 \\
7.90 & 1600 & 252.40
\end{tabular}
\end{center}

The spot price of the index corresponding to these option prices was $1,370$. For the
options above, the market estimates for the annualized continuous dividend yield of the
S\&P 500 index and for the risk-free rate were $q=0.0193=1.93\%$ and $r=0.0015=0.15\%$,
respectively.

\begin{subparts}
\item Use Newton's method to compute the implied volatilities for each of the options.
\item How do these values compare to the implied volatilities obtained by using
ordinary least squares to compute the discount factor and the present value of the
forward price?
\item How do the implied volatilities of calls and puts with the same strike compare to
each other?
\end{subparts}

\item On May 22, 2014, the bid and ask prices of the S\&P 500 options maturing on
January 17, 2015, were as follows:
\begin{center}
\begin{tabular}{ccccc}
Bid Price & Ask Price & Strike & Bid Price & Ask Price \\
Call & Call & & Put & Put \\[0.4em]
431.20 & 434.40 & 1450 & 8.90 & 10.20 \\
384.10 & 387.40 & 1500 & 11.60 & 13.10 \\
337.90 & 341.20 & 1550 & 15.20 & 16.90 \\
292.80 & 296.20 & 1600 & 19.90 & 21.80 \\
228.30 & 231.10 & 1675 & 29.70 & 31.40 \\
207.50 & 210.10 & 1700 & 33.60 & 35.50 \\
167.60 & 170.00 & 1750 & 43.50 & 45.40 \\
148.90 & 151.30 & 1775 & 49.40 & 51.50 \\
130.60 & 132.80 & 1800 & 56.50 & 58.10 \\
113.20 & 115.30 & 1825 & 63.60 & 65.70 \\
96.80 & 98.60 & 1850 & 71.90 & 74.20 \\
81.50 & 83.20 & 1875 & 81.40 & 83.50 \\
67.10 & 69.00 & 1900 & 92.10 & 94.50 \\
54.20 & 55.90 & 1925 & 103.80 & 106.30 \\
32.80 & 34.50 & 1975 & 132.40 & 134.70 \\
24.50 & 25.90 & 2000 & 149.20 & 151.90 \\
12.40 & 13.60 & 2050 & 186.90 & 189.70 \\
5.60 & 6.50 & 2100 & 230.00 & 232.40
\end{tabular}
\end{center}

The spot price of the index corresponding to these option prices was $1,894$.

\begin{subparts}
\item Compute the mid prices of the options, i.e., the average of the bid price and ask
price of the options. Use ordinary least squares to compute the present value of the
forward price $PVF$ and the discount factor $disc$ corresponding to the mid prices of
the options.
\item Compute the implied volatilities of these options. How do the implied
volatilities of calls and puts with the same strike compare to each other?
\end{subparts}

\item Recall that the vega of an option measures the sensitivity of the price of the
option to changes in volatility. In other words, for call and put options,
$\text{vega}(C) = \pd{C}{\sigma}$ and $\text{vega}(P) = \pd{P}{\sigma}$.

Use the Put--Call parity
% UNCLEAR/likely printing error, PDF page 262: printed with a plus sign
% before Ke^{-rT} (and a stray ink dot before "Se^{-qT}", almost certainly
% scan/print noise) -- verified at 1200 dpi. Standard Put-Call parity, and
% the identical relation restated in exercise 6 below, both give
% C - P = Se^{-qT} - Ke^{-rT} (minus). Transcribed here exactly as printed
% rather than silently corrected; see also PAGES 254-256 file's fidelity
% notes for this book's convention on printed errors.
\[
C - P \;=\; Se^{-qT} + Ke^{-rT}
\]
to show that
\[
\text{vega}(C) \;=\; \text{vega}(P).
\]

\item The goal of this exercise is to derive the Black--Scholes formulas for the vegas
of plain vanilla European calls and puts.

In the Black--Scholes framework, an underlying asset with spot price $S$ follows a
lognormal distribution with volatility $\sigma$ and pays continuous dividends at rate
$q$. Denote by $r$ be the risk-free interest rate, assumed to be constant. The
Black--Scholes values of a European call option and of a European put option with
strike $K$ and maturity $T$ are
\begin{align*}
C_{BS}(S,K,T,\sigma,r,q) &= Se^{-qT}N(d_1) - Ke^{-rT}N(d_2);\\
P_{BS}(S,K,T,\sigma,r,q) &= Ke^{-rT}N(-d_2) - Se^{-qT}N(-d_1),
\end{align*}
respectively, where
\[
N(z) \;=\; \frac{1}{\sqrt{2\pi}}\int_{-\infty}^z e^{-\frac{x^2}{2}}\,dx
\]
denotes the cumulative distribution of the standard normal variable, and
\[
d_1 \;=\; \frac{\ln\left(\frac{S}{K}\right) + \left(r-q+\frac{\sigma^2}{2}\right)T}
{\sigma\sqrt{T}}; \qquad d_2 \;=\; d_1 - \sigma\sqrt{T}.
\]

\begin{subparts}
\item Use Chain Rule to show that
\[
\text{vega}(C) \;=\; \pd{C}{\sigma} \;=\; Se^{-qT}N'(d_1)\pd{d_1}{\sigma} -
Ke^{-rT}N'(d_2)\pd{d_2}{\sigma}.
\]
\item Show that
\[
N'(z) \;=\; \frac{1}{\sqrt{2\pi}}e^{-\frac{z^2}{2}},
\]
and use this fact to prove that
\[
Se^{-qT}N'(d_1) \;=\; Ke^{-rT}N'(d_2).
\]
\item Show that
\begin{align*}
\text{vega}(C) &= \frac{1}{\sqrt{2\pi}}Se^{-qT}e^{-\frac{d_1^2}{2}}\sqrt{T};\\
\text{vega}(P) &= \frac{1}{\sqrt{2\pi}}Se^{-qT}e^{-\frac{d_1^2}{2}}\sqrt{T}.
\end{align*}
\end{subparts}

\item The goal of this exercise is to show that the theoretical values of the implied
volatilities of plain vanilla European call and put options with the same strike and
maturity are equal.

Denote by $C_{BS}(S,K,T,\sigma,r,q)$ and $P_{BS}(S,K,T,\sigma,r,q)$ the Black--Scholes
values of a European call option and of a European put option, respectively, with
strike $K$ and maturity $T$ on an underlying asset following a lognormal distribution
with volatility $\sigma$, with spot price $S$, and paying dividends continuously at the
rate $q$, if interest rates are constant and equal to $r$. Let $C_m$ and $P_m$ be the
market prices of a call and of a put option, respectively, with parameters $S$, $K$,
$T$, $r$, and $q$.

The implied volatility $\sigma_{imp,C}$ corresponding to the price $C$ is the solution
to
\[
C_{BS}(\sigma_{imp,C}) \;=\; C_m.
\]

The implied volatility $\sigma_{imp,P}$ corresponding to the price $P$ is the solution
to
\[
P_{BS}(\sigma_{imp,P}) \;=\; P_m.
\]

Note that, here and below, $C_{BS}(\sigma)$ and $P_{BS}(\sigma)$ are shorthand notations
for $C_{BS}(S,K,T,\sigma,r,q)$ and $P_{BS}(S,K,T,\sigma,r,q)$, respectively.

\begin{subparts}
\item Use the facts that the Black--Scholes values of put and call options satisfy the
Put--Call parity for any value $\sigma > 0$ of the volatility, i.e.,
\[
C_{BS}(\sigma) - P_{BS}(\sigma) \;=\; Se^{-qT} - Ke^{-rT},
\]
and that the market prices of put and call options satisfy the Put--Call parity as
well, i.e.,
\[
C_m - P_m \;=\; Se^{-qT} - Ke^{-rT},
\]
to conclude that
\[
P_{BS}(\sigma_{imp,P}) \;=\; P_{BS}(\sigma_{imp,C}).
\]
\item Show that the Black--Scholes value of a put option is a strictly increasing
function of volatility, and conclude that the implied volatilities corresponding to put
and call options with the same strike and maturity on the same asset must be equal,
i.e.,
\[
\sigma_{imp,P} \;=\; \sigma_{imp,C}.
\]
\end{subparts}

\item Recall the example from the book where for 15 consecutive trading days, the
yields of the 2-year, 3-year, 5-year, and 10-year treasury bonds were, respectively:
\begin{center}
\begin{tabular}{cccc}
2-year & 3-year & 5-year & 10-year \\[0.4em]
4.69 & 4.58 & 4.57 & 4.63 \\
4.81 & 4.71 & 4.69 & 4.73 \\
4.81 & 4.72 & 4.70 & 4.74 \\
4.79 & 4.78 & 4.77 & 4.81 \\
4.79 & 4.77 & 4.77 & 4.80 \\
4.83 & 4.75 & 4.73 & 4.79 \\
4.81 & 4.71 & 4.72 & 4.76 \\
4.81 & 4.72 & 4.74 & 4.77 \\
4.83 & 4.76 & 4.77 & 4.80 \\
4.81 & 4.73 & 4.75 & 4.77 \\
4.82 & 4.75 & 4.77 & 4.80 \\
4.82 & 4.75 & 4.76 & 4.80 \\
4.80 & 4.73 & 4.75 & 4.78 \\
4.78 & 4.71 & 4.72 & 4.73 \\
4.79 & 4.71 & 4.71 & 4.73
\end{tabular}
\end{center}

Denote by $T_2$, $T_3$, $T_5$, and $T_{10}$ the time series data vectors corresponding
to the yield of the 2-year, 3-year, 5-year, and 10-year treasury bonds, respectively.

\begin{subparts}
\item Find the coefficients $a$, $b_1$, $b_2$, $b_3$ of the linear regression for the
yield of the 3-year bond in terms of the yields of the 2-year, 5-year, and 10-year
bonds, i.e., find $a$, $b_1$, $b_2$, $b_3$ corresponding to the solution to the
ordinary least squares problem
\[
T_3 \;\approx\; a\bfone + b_1T_2 + b_2T_5 + b_3T_{10},
\]
where $\bfone$ is the $15\times1$ column vector with all entries equal to $1$. Let
\[
T_{3,LR} \;=\; a\bfone + b_1T_2 + b_2T_5 + b_3T_{10}.
\]
Find the approximation error
\[
\text{error}_{LR} \;=\; \norm{T_3 - T_{3,LR}}
\]
of the linear regression.

\item Compute the linear interpolation values of the 3-year yield by doing linear
interpolation between the 2-year yield and the 5-year yield at each data point. Denote
by $T_{3,linear\_interp}$ the time series vector of these values. In other words,
\[
T_{3,linear\_interp} \;=\; \frac23T_2 + \frac13T_5.
\]
Find the approximation error
\[
\text{error}_{linear\_interp} \;=\; \norm{T_3 - T_{3,linear\_interp}}
\]
of the linear interpolation.

\item Compute the cubic interpolation values of the 3-year yield by doing cubic spline
interpolation between the 2-year, 5-year, and 10-year yield at each data point. Denote
by $T_{3,cubic\_interp}$ the time series vector of these values.

Find the approximation error
\[
\text{error}_{cubic\_interp} \;=\; \norm{T_3 - T_{3,cubic\_interp}}
\]
of the cubic interpolation.

\item Compare the approximation errors from (i), (ii), and (iii), and comment on the
results.
\end{subparts}

\item The file \emph{financials2012.xlsx} from www.fepress.org/nla-primer contains the
end of week adjusted closing prices for the stocks of the following financial
companies: JPM; GS; MS; BAC (Bank of America); RBS; CS; UBS; RY (RBC); BCS (Barclays)
between January 11, 2012, and October 15, 2012.

\begin{subparts}
\item Compute the weekly percentage returns of these stocks.
\item Find the linear regression of the JPM returns with respect to the returns of the
other stocks. What is the approximation error of this linear regression?
\item Find the linear regression of the JPM returns with respect to the returns of the
other American financial companies, i.e., with respect to GS, MS, and BAC. What is the
approximation error of this linear regression?
\item Find the linear regression of the JPM stock prices with respect to the prices of
the other stocks. What is the approximation error of this linear regression? How does
it compare with the approximation error of the linear regression of the JPM returns
computed at (ii)?
\end{subparts}

\end{exercises}
